\chapter*{符号表}
\addcontentsline{toc}{chapter}{符号表}
\markboth{符号表}{符号表}

原书页码沿用英文版符号表；译本定位链接到相应章的起页。

\begingroup
\small\raggedright
\begin{longtable}{@{}p{.26\linewidth}p{.41\linewidth}p{.12\linewidth}p{.12\linewidth}@{}}
\toprule 符号 & 说明 & 原书页码 & 译本定位\\\midrule
\endfirsthead
\toprule 符号 & 说明 & 原书页码 & 译本定位\\\midrule
\endhead
\(\clubsuit_n\)（原书 \(\maltese_n\)） & 符号指数 \(|a_1|+\cdots+|a_n|-n\) & 8 & \pageref{ch:01-definitions}\\
\(\operatorname{Aut}(M)\) & 正合辛自同构群 & 96 & \pageref{ch:07-symplectic-geometry}\\
\(\operatorname{Aut}(\widehat M)\) & 接上锥之后的辛自同构群 & 98 & \pageref{ch:07-symplectic-geometry}\\
\(\operatorname{Aut}(M,\partial M)\) & 相对于边界的正合辛自同构群 & 96 & \pageref{ch:07-symplectic-geometry}\\
\(\operatorname{Aut}^{\mathrm{gr}}(M,\partial M)\) & 分次辛自同构群 & 188 & \pageref{ch:12-fukaya-complete}\\
\(\operatorname{Aut}^{\#}(M,\partial M)\) & 膜自同构群 & 188 & \pageref{ch:12-fukaya-complete}\\
\(CF^{\mathrm{pr}}(L_0,L_1)\) & 弗洛尔余链群（初步版本） & 109 & \pageref{ch:08-classical-floer}\\
\(CF^*(L_0^\#,L_1^\#)\) & 弗洛尔余链群 & 181 & \pageref{ch:12-fukaya-complete}\\
\(\operatorname{Crit}(\pi)\) & 临界点集 & 212 & \pageref{ch:15-first-notions}\\
\(\operatorname{Critv}(\pi)\) & 临界值集 & 212 & \pageref{ch:15-first-notions}\\
\(D(\mathcal A),D^\pi(\mathcal A)\) & 导出范畴 & 42 & \pageref{ch:03-exact-triangles}\\
\(D_{S,u}\) & 线性化算子 & 107 & \pageref{ch:08-classical-floer}\\
\(D_{\mathcal S,r,u}\) & 扩大的线性化算子 & 122 & \pageref{ch:09-fukaya-preliminary}\\
\(\partial^h E\) & 水平边界 & 212 & \pageref{ch:15-first-notions}\\
\(\partial^v E\) & 竖直边界 & 212 & \pageref{ch:15-first-notions}\\
\(\operatorname{Fakev}(\Psi,\varpi)\) & 伪临界值集 & 220 & \pageref{ch:15-first-notions}\\
\(\mathcal F(M)^{\mathrm{pr}}\) & 深谷范畴（初步版本） & 126 & \pageref{ch:09-fukaya-preliminary}\\
\(\mathcal F(\pi)\) & 莱夫谢茨纤维化的深谷范畴 & 282 & \pageref{ch:18-fukaya-lefschetz}\\
\(H(\mathcal A),H^0(\mathcal A)\) & 上同调范畴 & 8 & \pageref{ch:01-definitions}\\
\(H\mathcal F(M)\) & 唐纳森--深谷范畴 & 181 & \pageref{ch:12-fukaya-complete}\\
\(H\mathcal F(M)^{\mathrm{pr}}\) & 唐纳森--深谷范畴（初步版本） & 101 & \pageref{ch:08-classical-floer}\\
\(HF^{\mathrm{pr}}\) & 弗洛尔上同调（初步版本） & 101, 109 & \pageref{ch:08-classical-floer}\\
\(HF^*(L_0^\#,L_1^\#)\) & 弗洛尔上同调 & 182 & \pageref{ch:12-fukaya-complete}\\
\(\mathcal M^{d+1}(y_0,\ldots,y_d)\) & 伪全纯多边形模空间 & 126 & \pageref{ch:09-fukaya-preliminary}\\
\(\mathcal M_{E/S}(\{y_\xi\})\) & 非齐次伪全纯截面的模空间 & 240 & \pageref{ch:17-pseudoholomorphic-sections}\\
\(\mathcal M_S(\{y_\xi\}),\mathcal M_{\mathcal S}(\{y_\xi\})\) & 伪全纯映射的模空间 & 105, 122 & \pageref{ch:09-fukaya-preliminary}\\
\(\mathcal M_Z^*(y_0,y_1)\) & 弗洛尔轨线的模空间 & 105 & \pageref{ch:08-classical-floer}\\
\(\operatorname{nu\text{-}fun}(\mathcal A,\mathcal B),\operatorname{fun}(\mathcal A,\mathcal B)\) & \(A_\infty\)-函子范畴 & 10, 28 & \pageref{ch:02-identity-morphisms}\\
\(\operatorname{nu\text{-}mod}(\mathcal A),\operatorname{mod}(\mathcal A)\) & \(A_\infty\)-模范畴 & 19, 30 & \pageref{ch:02-identity-morphisms}\\
\(\mathcal R^{d+1},\overline{\mathcal R}^{d+1}\) & 带标记圆盘的模空间及其紧化 & 115, 118 & \pageref{ch:09-fukaya-preliminary}\\
\(TE^v,TE^h\) & 沿纤维的切空间及其补空间 & 212 & \pageref{ch:15-first-notions}\\
\(T_{Y_0}(Y_1)\) & 对象的扭转 & 63 & \pageref{ch:04-idempotents}\\
\(Z\) & 无穷带状域 & 100 & \pageref{ch:08-classical-floer}\\
\(Z_\pm\) & 半无穷带状域 & 103 & \pageref{ch:08-classical-floer}\\
\bottomrule
\end{longtable}
\endgroup
